% !PRINTED none
\begin{center}
\gesperrt{Seventh Section.}

\medskip
{\Large\textbf{Theory of the Transformation Equations.}}

\medskip
\rule{3cm}{0.4pt}
\end{center}

\begin{center}
\textbf{\S\,66. \gesperrt{Formation of Transformation Equations.}}
\end{center}

Now that the division equations have been reduced to transformation equations, we proceed to a more exact study of these latter equations.

We assume $n$ to be odd, understand by $p$ the primes that divide $n$, and put

\weq{1}{\nu = \psi(n) = n\,\Pi\left(1+\frac{1}{p}\right),}

and denote the $\nu$ rows of the roots of the division equation by $R_1, R_2, \ldots, R_\nu$.

From each of these rows we take, for

\weq{2}{\Omega_{\mu,\mu'} = \frac{4\mu K + 4\mu' i K'}{n}}

a representative $\Omega_1, \Omega_2, \ldots, \Omega_\nu$, and we obtain the $\varphi(n)$ roots of one row if in

\[
\mathrm{sn}\, m\,\Omega
\]

we let $m$ run through a complete system of incongruent numbers coprime to $n$.

The simplest expressions that can be introduced as roots of transformation equations are the products

% [sic] The source prints n above the product sign, although the following sentence says that h is the running index.
\weq{3}{\mathop{\Pi}\limits_{1,\,n-1}^{n} \Phi(\mathrm{sn}\,h\,\Omega),}

where $\Phi$ is an arbitrary rational function, and $h$ runs through the sequence of numbers $1, 2, \ldots, n-1$.

Every such function is rationally expressible by $\mathrm{sn}\,\Omega$, and evidently remains unchanged when $\Omega$ is replaced by any $m\,\Omega$;
%CONT>>
